% GENERATED FILE. Edit canonical lessons, then run npm run generate:books.
% canonical-source-hash: fnv1a64:d08010f162b51d8f
% canonical-lessons: FR-C360-augmentation, FR-C360-promotion, FR-C360-qualification, FR-C360-competence, FR-C360-formation

\chapter{Faire carrière}
\label{ch:fr-faire-carriere}
\hlchaptermodality{360}

\begin{chapteropening}
\textbf{By the end of this chapter:} \emph{I can talk about getting on at work: a pay rise, a promotion, qualifications, skills and training.}

Complete the last lesson of chapter 360: I can talk about getting on at work: a pay rise, a promotion, qualifications, skills and training.
\end{chapteropening}

\section[l'augmentation]{l'augmentation — the pay rise, the increase}

\label{lesson:FR-C360-augmentation}



\begin{quote}
Before the new one: say the French for the extension, then the French for the delivery.
\end{quote}

\subsection*{You'll want to know: l'augmentation}
\textbf{l'augmentation} — \textquotedblleft{}the pay rise, the increase\textquotedblright{}.

\begin{quote}
J'ai demandé une augmentation à mon patron.
\end{quote}

\begin{quote}
\emph{I asked my boss for a pay rise.}
\end{quote}

From \emph{augmenter}. Feminine.

\begin{cousinweb}
\emph{augmentation} comes from Late Latin \emph{augmentātiō} \textquotedblleft{}an increase\textquotedblright{} (EtymDB), from \emph{augēre} \textquotedblleft{}to increase\textquotedblright{} (Lewis \& Short).
\end{cousinweb}

\begin{grammarlens}[title={What you've built}]
One more word for telling, explaining and describing at B1.
\end{grammarlens}

\subsection*{Your turn}
\begin{itemize}
\raggedright
  \item \emph{Say it:} \emph{l'augmentation}
  \item \emph{Say it:} \emph{l'augmentation}, once more
  \item \emph{Say it:} a sentence of your own with \emph{l'augmentation}
  \item \emph{Recall:} say the French for the extension, then the French for the delivery, then say \emph{l'augmentation} again
\end{itemize}

\begin{tcolorbox}[breakable,colback=teal!4,colframe=teal!35!black,title={Before you move on}]
What does l'augmentation mean? (\textquotedblleft{}The pay rise, the increase\textquotedblright{}.) Say it once more.
\end{tcolorbox}
\section[la promotion]{la promotion — the promotion; the special offer}

\label{lesson:FR-C360-promotion}



\begin{quote}
Before the new one: say the French for the delivery, then the French for the pay rise.
\end{quote}

\subsection*{You'll want to know: la promotion}
\textbf{la promotion} — \textquotedblleft{}the promotion; the special offer\textquotedblright{}.

\begin{quote}
Elle a eu une promotion : elle dirige maintenant l'équipe.
\end{quote}

\begin{quote}
\emph{She got a promotion: she now runs the team.}
\end{quote}

In shops, \emph{en promotion} — on special offer.

\begin{cousinweb}
\emph{promotion} is a learned borrowing from Latin \emph{prōmōtiō} \textquotedblleft{}advancement\textquotedblright{} (EtymDB), from \emph{prōmovēre} \textquotedblleft{}to move forward\textquotedblright{}, built on \emph{movēre} \textquotedblleft{}to move\textquotedblright{} (Lewis \& Short).
\end{cousinweb}

\begin{grammarlens}[title={What you've built}]
One more word for telling, explaining and describing at B1.
\end{grammarlens}

\subsection*{Your turn}
\begin{itemize}
\raggedright
  \item \emph{Say it:} \emph{la promotion}
  \item \emph{Say it:} \emph{la promotion}, once more
  \item \emph{Say it:} a sentence of your own with \emph{la promotion}
  \item \emph{Recall:} say the French for the delivery, then the French for the pay rise, then say \emph{la promotion} again
\end{itemize}

\begin{tcolorbox}[breakable,colback=teal!4,colframe=teal!35!black,title={Before you move on}]
What does la promotion mean? (\textquotedblleft{}The promotion; the special offer\textquotedblright{}.) Say it once more.
\end{tcolorbox}
\section[la qualification]{la qualification — the qualification}

\label{lesson:FR-C360-qualification}



\begin{quote}
Before the new one: say the French for the pay rise, then the French for the promotion.
\end{quote}

\subsection*{You'll want to know: la qualification}
\textbf{la qualification} — \textquotedblleft{}the qualification\textquotedblright{}.

\begin{quote}
Ce travail ne demande pas de qualification.
\end{quote}

\begin{quote}
\emph{This job does not require any qualification.}
\end{quote}

In sport, \emph{la qualification} is qualifying for the next round.

\begin{cousinweb}
\emph{qualification} comes from Medieval Latin \emph{quālificātiō}, built on Latin \emph{quālis} \textquotedblleft{}of what kind\textquotedblright{} (Lewis \& Short).
\end{cousinweb}

\begin{grammarlens}[title={What you've built}]
One more word for telling, explaining and describing at B1.
\end{grammarlens}

\subsection*{Your turn}
\begin{itemize}
\raggedright
  \item \emph{Say it:} \emph{la qualification}
  \item \emph{Say it:} \emph{la qualification}, once more
  \item \emph{Say it:} a sentence of your own with \emph{la qualification}
  \item \emph{Recall:} say the French for the pay rise, then the French for the promotion, then say \emph{la qualification} again
\end{itemize}

\begin{tcolorbox}[breakable,colback=teal!4,colframe=teal!35!black,title={Before you move on}]
What does la qualification mean? (\textquotedblleft{}The qualification\textquotedblright{}.) Say it once more.
\end{tcolorbox}
\section[la compétence]{la compétence — the skill, the competence}

\label{lesson:FR-C360-competence}



\begin{quote}
Before the new one: say the French for the promotion, then the French for the qualification.
\end{quote}

\subsection*{You'll want to know: la compétence}
\textbf{la compétence} — \textquotedblleft{}the skill, the competence\textquotedblright{}.

\begin{quote}
Il a de bonnes compétences en langues.
\end{quote}

\begin{quote}
\emph{He has good language skills.}
\end{quote}

Mostly in the plural for skills: \emph{des compétences}.

\begin{cousinweb}
\emph{compétence} comes from Late Latin \emph{competentia} (EtymDB), from \emph{competere} \textquotedblleft{}to come together, to be suitable\textquotedblright{}, built on \emph{petere} \textquotedblleft{}to seek\textquotedblright{} (Lewis \& Short).
\end{cousinweb}

\begin{grammarlens}[title={What you've built}]
One more word for telling, explaining and describing at B1.
\end{grammarlens}

\subsection*{Your turn}
\begin{itemize}
\raggedright
  \item \emph{Say it:} \emph{la compétence}
  \item \emph{Say it:} \emph{la compétence}, once more
  \item \emph{Say it:} a sentence of your own with \emph{la compétence}
  \item \emph{Recall:} say the French for the promotion, then the French for the qualification, then say \emph{la compétence} again
\end{itemize}

\begin{tcolorbox}[breakable,colback=teal!4,colframe=teal!35!black,title={Before you move on}]
What does la compétence mean? (\textquotedblleft{}The skill, the competence\textquotedblright{}.) Say it once more.
\end{tcolorbox}
\section[la formation]{la formation — the training, the education}

\label{lesson:FR-C360-formation}



\begin{quote}
Before the new one: say the French for the qualification, then the French for the skill.
\end{quote}

\subsection*{You'll want to know: la formation}
\textbf{la formation} — \textquotedblleft{}the training, the education\textquotedblright{}.

\begin{quote}
L'entreprise paie ma formation.
\end{quote}

\begin{quote}
\emph{The company is paying for my training.}
\end{quote}

\emph{La formation continue} — training for adults in work.

\begin{cousinweb}
\emph{formation} comes through Old French \emph{formacion} from Latin \emph{fōrmātiō} \textquotedblleft{}a shaping\textquotedblright{} (EtymDB), from \emph{fōrmāre} \textquotedblleft{}to shape\textquotedblright{}, built on \emph{fōrma} \textquotedblleft{}shape\textquotedblright{} (Lewis \& Short).
\end{cousinweb}

\begin{grammarlens}[title={What you've built}]
One more word for telling, explaining and describing at B1.
\end{grammarlens}

\subsection*{Your turn}
\begin{itemize}
\raggedright
  \item \emph{Say it:} \emph{la formation}
  \item \emph{Say it:} \emph{la formation}, once more
  \item \emph{Say it:} a sentence of your own with \emph{la formation}
  \item \emph{Recall:} say the French for the qualification, then the French for the skill, then say \emph{la formation} again
\end{itemize}

\begin{tcolorbox}[breakable,colback=teal!4,colframe=teal!35!black,title={Before you move on}]
What does la formation mean? (\textquotedblleft{}The training, the education\textquotedblright{}.) Say it once more.
\end{tcolorbox}
